\section{Rezultati}

\subsection{Struktura ispitanika}

Struktura ispitanika prikazana je kroz osnovne demografske raspodjele. Detaljne tablice generiraju se iz analitičke skripte i nalaze se u direktoriju \texttt{report/tables/}.

\begin{table}[H]
\centering
\caption{Raspodjela ispitanika prema spolu}
\begin{tabular}{lrr}
\toprule
Odgovor & Broj & Postotak \\
\midrule
Muški & 520 & 54.51 \\
Ženski & 434 & 45.49 \\
\bottomrule
\end{tabular}

\end{table}

\begin{table}[H]
\centering
\caption{Raspodjela ispitanika prema dobi}
\begin{tabular}{lrr}
\toprule
Odgovor & Broj & Postotak \\
\midrule
15-24 godina & 637 & 66.49 \\
40-54 godina & 144 & 15.03 \\
25-39 godina & 89 & 9.29 \\
55 i više godina & 88 & 9.19 \\
\bottomrule
\end{tabular}

\end{table}

\subsection{P9: najveća percipirana prednost SMR tehnologije}

Pitanje P9 ispituje što ispitanici smatraju najvećom prednošću SMR tehnologije. Rezultati su prikazani u tablici \ref{tab:p9} i na slici \ref{fig:p9}.

\begin{table}[H]
\centering
\caption{Frekvencije odgovora na pitanje P9}
\label{tab:p9}
\begin{tabular}{lrr}
\toprule
Odgovor & Broj & Postotak \\
\midrule
Niži troškovi & 154 & 16.19 \\
Kraća izgradnja & 277 & 29.13 \\
Sigurnosne značajke & 322 & 33.86 \\
Ne znam & 198 & 20.82 \\
\bottomrule
\end{tabular}

\end{table}

\begin{figure}[H]
\centering
\includegraphics[width=0.9\textwidth]{p9_odgovori.pdf}
\caption{Raspodjela odgovora na pitanje P9}
\label{fig:p9}
\end{figure}

\subsection{P10: stav prema državnim subvencijama}

Pitanje P10 ispituje smatraju li ispitanici opravdanim državnu subvenciju razvoja SMR tehnologije umjesto ulaganja u jednu klasičnu veliku nuklearnu elektranu. Rezultati su prikazani u tablici \ref{tab:p10} i na slici \ref{fig:p10}.

\begin{table}[H]
\centering
\caption{Frekvencije odgovora na pitanje P10}
\label{tab:p10}
\begin{tabular}{lrr}
\toprule
Odgovor & Broj & Postotak \\
\midrule
Da & 560 & 59.01 \\
Ne & 76 & 8.01 \\
Ne znam & 313 & 32.98 \\
\bottomrule
\end{tabular}

\end{table}

\begin{figure}[H]
\centering
\includegraphics[width=0.9\textwidth]{p10_odgovori.pdf}
\caption{Raspodjela odgovora na pitanje P10}
\label{fig:p10}
\end{figure}

\subsection{Demografske usporedbe i statistički testovi}

Za pitanja P9 i P10 provedene su usporedbe prema spolu, dobi, obrazovanju, zaposlenosti i županiji. Sažetak hi-kvadrat testova prikazan je u tablici \ref{tab:testovi}. U konačnoj interpretaciji naglasak treba staviti na statistički značajne i sadržajno relevantne povezanosti.

Prije provedbe hi-kvadrat testova kategorije s manje od pet opažanja spojene su u zajedničku kategoriju \textit{Ostalo (<5)}. Time se smanjuje utjecaj vrlo rijetkih kategorija na testnu statistiku i dobiva stabilnija procjena povezanosti.

\begin{table}[H]
\centering
\caption{Rezultati hi-kvadrat testova neovisnosti}
\label{tab:testovi}
\resizebox{\textwidth}{!}{\begin{tabular}{llrrrrrl}
\toprule
Varijabla 1 & Varijabla 2 & Hi-kvadrat & df & p-vrijednost & Cramér V & Najmanja očekivana frekvencija & Napomena \\
\midrule
Županija & P9 & 101.13 & 90 & 0.20 & 0.19 & 0.79 & Oprez: dio očekivanih frekvencija je manji od 5. \\
Županija & P10 & 56.52 & 60 & 0.60 & 0.18 & 0.38 & Oprez: dio očekivanih frekvencija je manji od 5. \\
Spol & P9 & 75.89 & 3 & 0.00 & 0.28 & 69.09 &  \\
Spol & P10 & 5.79 & 2 & 0.06 & 0.08 & 34.16 &  \\
Dob & P9 & 13.11 & 9 & 0.16 & 0.07 & 14.10 &  \\
Dob & P10 & 7.15 & 6 & 0.31 & 0.06 & 7.05 &  \\
Obrazovanje & P9 & 11.83 & 12 & 0.46 & 0.06 & 0.81 & Oprez: dio očekivanih frekvencija je manji od 5. \\
Obrazovanje & P10 & 19.11 & 8 & 0.01 & 0.10 & 0.40 & Oprez: dio očekivanih frekvencija je manji od 5. \\
Status zaposlenosti & P9 & 15.10 & 9 & 0.09 & 0.07 & 4.38 & Oprez: dio očekivanih frekvencija je manji od 5. \\
Status zaposlenosti & P10 & 5.22 & 6 & 0.52 & 0.05 & 2.22 & Oprez: dio očekivanih frekvencija je manji od 5. \\
P9 & P10 & 191.38 & 6 & 0.00 & 0.32 & 12.20 &  \\
\bottomrule
\end{tabular}
}
\end{table}

\subsection{Grafički prikaz statistički značajnih povezanosti}

% Automatski generirano skriptom analysis/src/analyze.py
U nastavku su prikazani grafovi povezanosti za parove varijabli kod kojih je hi-kvadrat test pokazao statistički značajnu povezanost na razini značajnosti $\alpha = 0{,}05$. Prije testiranja kategorije s manje od pet opažanja spojene su u kategoriju \textit{Ostalo (<5)}.

\begin{figure}[H]
\centering
\includegraphics[width=0.92\textwidth]{povezanost_spol_p9_kratko.pdf}
\caption{Povezanost varijabli Spol i P9}
\label{fig:povezanost_spol_p9_kratko}
\end{figure}

\begin{figure}[H]
\centering
\includegraphics[width=0.92\textwidth]{povezanost_obrazovanje_p10_kratko.pdf}
\caption{Povezanost varijabli Obrazovanje i P10}
\label{fig:povezanost_obrazovanje_p10_kratko}
\end{figure}

\begin{figure}[H]
\centering
\includegraphics[width=0.92\textwidth]{povezanost_p9_kratko_p10_kratko.pdf}
\caption{Povezanost varijabli P9 i P10}
\label{fig:povezanost_p9_kratko_p10_kratko}
\end{figure}


\subsection{Povezanost odgovora na P9 i P10}

Križna tablica odgovora na P9 i P10 prikazuje jesu li stavovi o glavnoj prednosti SMR tehnologije povezani sa stavom prema državnom subvencioniranju razvoja te tehnologije.

\begin{table}[H]
\centering
\caption{Križna tablica odgovora P9 i P10}
\resizebox{\textwidth}{!}{\begin{tabular}{lrrrr}
\toprule
P10 & Da & Ne & Ne znam & Ukupno \\
P9 &  &  &  &  \\
\midrule
Kraća izgradnja & 188 & 25 & 63 & 276 \\
Ne znam & 38 & 17 & 141 & 196 \\
Niži troškovi & 94 & 17 & 43 & 154 \\
Sigurnosne značajke & 240 & 16 & 65 & 321 \\
Ukupno & 560 & 75 & 312 & 947 \\
\bottomrule
\end{tabular}
}
\end{table}
